\section*{Question 3}

We model the normalized expression $Y$ for each gene by the simple linear
regression
\[
    Y = b_0 + b_1 x + e ,
\]
where $x$ is an indicator of condition:
\[
    x =
    \begin{cases}
        0, & \text{control group C},   \\
        1, & \text{treatment group T}.
    \end{cases}
\]
For a binary predictor with equal numbers of samples in the two groups,
$b_0$ and $b_1$ have the interpretation
\[
    b_0 = \text{mean expression in group C}, \qquad
    b_1 = \text{mean}(Y\mid x=1) - \text{mean}(Y\mid x=0),
\]
i.e.\ $b_1$ is the difference between the treatment and control means.

%---------------------------------------------------------------
\subsubsection*{1. Estimation of $b_0$ and $b_1$}

For each gene we compute
\[
    \bar{T} = \frac{T_1+T_2+T_3}{3}, \qquad
    \bar{C} = \frac{C_1+C_2+C_3}{3},
\]
then set $b_0 = \bar{C}$ and $b_1 = \bar{T}-\bar{C}$.

As an example, for Gene 1
\[
    \bar{T}_1 = \frac{4.1+5.2+2.8}{3} = 4.03\overline{3}, \qquad
    \bar{C}_1 = \frac{5.6+3.7+3.8}{3} = 4.36\overline{6},
\]
so
\[
    b_{0,1} = 4.36\overline{6}, \qquad
    b_{1,1} = 4.03\overline{3} - 4.36\overline{6} = -0.33\overline{3}.
\]

Carrying out the same calculation for all six genes yields

\[
    \begin{array}{c|ccc}
        \text{Gene}   & \bar{C} & \bar{T} & b_1 = \bar{T}-\bar{C} \\ \hline
        \text{Gene 1} & 4.37    & 4.03    & -0.33                 \\
        \text{Gene 2} & 4.30    & 2.73    & -1.57                 \\
        \text{Gene 3} & 4.50    & 2.93    & -1.57                 \\
        \text{Gene 4} & 1.37    & 2.30    & 0.93                  \\
        \text{Gene 5} & 3.07    & 2.57    & -0.50                 \\
        \text{Gene 6} & 2.40    & 5.43    & 3.03
    \end{array}
\]

(Values are rounded to two decimal places.)

Thus the fitted regression models are
\[
    Y_g = b_{0,g} + b_{1,g}\,x + e_g, \qquad g=1,\dots,6.
\]

%---------------------------------------------------------------
\subsubsection*{2. Testing for differential expression}

We are told that, over all genes in the genome, the regression
coefficients $b_1$ follow an approximately normal distribution
(as illustrated by the given bell curve).  We treat the six
estimated coefficients $b_{1,g}$ above as a sample from
\[
    b_1 \sim N(0,\sigma^2).
\]

First we estimate $\sigma$ using the sample standard deviation of the six
coefficients:
\[
    \hat{\sigma}
    = \sqrt{\frac{1}{n-1}\sum_{g=1}^n (b_{1,g}-\bar{b}_1)^2}
    \approx 1.75,
\]
where $\bar{b}_1$ is the sample mean of the six $b_{1,g}$
(which is essentially $0$).

For each gene we then compute the standardized statistic
\[
    Z_g = \frac{b_{1,g}-\bar{b}_1}{\hat{\sigma}}
    \approx \frac{b_{1,g}}{1.75}.
\]

The resulting $Z$-scores are

\[
    \begin{array}{c|cc}
        \text{Gene}   & b_1   & Z_g \approx b_1/1.75 \\ \hline
        \text{Gene 1} & -0.33 & -0.19                \\
        \text{Gene 2} & -1.57 & -0.89                \\
        \text{Gene 3} & -1.57 & -0.89                \\
        \text{Gene 4} & 0.93  & 0.53                 \\
        \text{Gene 5} & -0.50 & -0.29                \\
        \text{Gene 6} & 3.03  & 1.73
    \end{array}
\]

We perform a two-sided test of
\[
    H_0: b_1 = 0 \quad \text{vs.} \quad H_A: b_1 \neq 0
\]
at significance level $\alpha = 0.10$.
For a standard normal distribution,
\[
    z_{0.95} \approx 1.645,
\]
so the rejection region is
\[
    |Z_g| > 1.645.
\]

Comparing the table above with this critical value, only Gene 6 has
\[
    |Z_6| \approx 1.73 > 1.645,
\]
while all other genes satisfy $|Z_g| < 1.645$.

%---------------------------------------------------------------
\subsubsection*{3. Conclusion}

At the $10\%$ level of significance, only \textbf{Gene 6} shows a
regression coefficient $b_1$ that lies in the outer $10\%$ of the
overall $b_1$ distribution. Therefore,
\[
    \boxed{\text{Gene 6 is the only differentially expressed gene.}}
\]
All other genes are not significantly differentially expressed between
conditions T and C at $\alpha = 0.1$.
